\subsection{Congruence bounds with repeated residual blocks}
\label{sec:residual-block-fitting}

We refine the elementary Fitting inequality of
Adriazola~\cite[\S17, Proposition \texttt{w4:prop:fitting}]{adriazola-reconstruction}.
Joint residual surjectivity removes its residue-field-size condition.
Individual surjectivity permits repeated residual constituents, with
the precise alternatives below. These algebraic statements do not
replace the arithmetic construction of the cocycle module or its
local estimates.

Let $O$ be a discrete valuation ring, $K$ its fraction field,
$\varpi$ a uniformizer, $k$ its residue field, and $v(\varpi)=1$.
Let $X$ be finite, $x_0\in X$, and $T\subset O^X$ an $O$-subalgebra.
Put
\[
 \mathfrak c_{x_0}=\{t(x_0):t\in T,\ t(x)=0\text{ for }x\ne x_0\},
 \qquad \delta_{x_0}=v(\mathfrak c_{x_0}),
\]
with $v(0)=+\infty$ for the zero ideal. Let $R$ be an $O$-algebra,
possibly noncommutative, and let
$\rho_x:R\to M_{d_x}(O)$ be $O$-algebra representations.
Fix $P\in M_{a\times b}(R)$ with $a\ge b\ge1$. Assume:
\begin{enumerate}
\item for every $Q\in M_b(R)$, the function
      $x\mapsto\det\rho_x(Q)$ belongs to $T$;
\item for $x\ne x_0$, the matrix $\rho_x(P)$ has a nonzero
      kernel over $K$;
\item in an integral basis $\rho_{x_0}$ is block upper triangular,
      with diagonal representations $\rho_j:R\to M_{n_j}(O)$,
      $1\le j\le g$, where $g,n_j\ge1$.
\end{enumerate}
Write $P_j=\rho_j(P)$, $D_0=b\sum_jn_j$, and
\[
 L_j=\operatorname{length}_O
 \ker\bigl(P_j:(K/O)^{bn_j}\to(K/O)^{an_j}\bigr).
\]
When all $L_j$ are finite, $P_j$ is injective over $K$.
Its saturated image $W_j\subset O^{an_j}$ is a free direct summand,
of rank $bn_j$, and elementary divisors give
\begin{equation}\label{eq:residual-block-length}
 L_j=\operatorname{length}_O(W_j/\operatorname{im}P_j).
\end{equation}

\begin{proposition}[Exact bound with joint residual image]
\label{prop:residual-block-joint}
If $R\to\prod_jM_{n_j}(k)$ is onto, then
\begin{equation}\label{eq:residual-block-exact}
                       \delta_{x_0}\le\sum_jL_j.
\end{equation}
There is no restriction on the size or characteristic of $k$.
\end{proposition}
\begin{proof}
If a length is infinite there is nothing to prove. Otherwise choose
for each $j$ a linear map
$Y_j:k^{an_j}\to k^{bn_j}$ that is an isomorphism on
$W_j/\varpi W_j$. Extend a basis of that subspace to obtain such
a map. Write $Y_j$ in $b\times a$ blocks of size $n_j$.
The joint surjection lifts the tuples of blocks, one position at a
time, to a matrix $C\in M_{b\times a}(R)$.
Consequently $\rho_j(C)|_{W_j}$ is an isomorphism over $O$.

Set $Q=CP$ and $f(x)=\det\rho_x(Q)$. Assumption 1 puts $f$ in $T$;
assumption 2 makes it vanish outside $x_0$. Block triangularity and
\eqref{eq:residual-block-length} give
\[
 v(f(x_0))=\sum_jv(\det\rho_j(CP))=\sum_jL_j.
\]
Thus $f(x_0)\in\mathfrak c_{x_0}$ proves the inequality.
\end{proof}

\begin{proposition}[Individual residual images]
\label{prop:residual-block-large-field}
Suppose each map $\bar\rho_j:R\to M_{n_j}(k)$ is onto.
If $k$ is infinite or $|k|>D_0$, then
\eqref{eq:residual-block-exact} holds.
\end{proposition}
\begin{proof}
Let $S\subset\prod_jM_{n_j}(k)$ be the joint image, viewed as a
finite-dimensional $k$-vector space. Choose a basis and lifts in $R$.
On $S^{b\times a}$, the determinant of the $j$th matrix restricted
to $W_j/\varpi W_j$ is a polynomial of degree $bn_j$.
It is nonzero: individual surjectivity permits a left inverse
on this subspace. The product of these determinant polynomials
is nonzero, since the coordinate polynomial ring of a vector
space is a domain, and its total degree is $D_0$.
A nonzero polynomial of total degree less than $|k|$ has a
nonzero value on $k$-points, by induction on the number of variables;
the assertion also holds over an infinite field.
Lift such a point into $R$. The resulting $C$ satisfies the
same invertibility conditions used in
Proposition~\ref{prop:residual-block-joint}.
\end{proof}

\begin{lemma}[Integral grid bound]
\label{lem:residual-block-grid}
Suppose $O$ has characteristic zero. Let
$F\in O[Z_1,\ldots,Z_t]$ have total degree at most $D$ and
at least one unit coefficient. There exist ordinary integers
$0\le z_i\le D$ such that
\[
                            v(F(z_1,\ldots,z_t))\le v(D!).
\]
\end{lemma}
\begin{proof}
The Newton expansion of $F$ uses the products
$\prod_i\binom{Z_i}{a_i}$, $\sum_i a_i\le D$.
Their coefficients are the iterated finite differences
$\Delta^aF(0)$, integral linear combinations of values on the stated
grid. Each denominator in the ordinary monomial expansion of such
a product divides $\prod_i a_i!$, and this integer divides $D!$.
If every grid value had valuation greater than $v(D!)$, all finite
differences would too. Conversion to ordinary monomials would then
put every coefficient of $F$ in the maximal ideal, a contradiction.
\end{proof}

\begin{theorem}[Repeated blocks with a factorial error]
\label{thm:residual-block-factorial}
Suppose $O$ has characteristic zero and residue characteristic
$\ell$, and each $\bar\rho_j$ is onto. Then
\begin{equation}\label{eq:residual-block-factorial}
                   \delta_{x_0}\le\sum_jL_j+v(D_0!).
\end{equation}
With $e_O=v(\ell)$, this becomes
\[
 \frac{\delta_{x_0}}{e_O}
       \le\sum_j\frac{L_j}{e_O}+v_\ell(D_0!),
 \qquad
 v_\ell(D_0!)=\sum_{a\ge1}\left\lfloor\frac{D_0}{\ell^a}\right\rfloor
       \le\frac{D_0}{\ell-1}.
\]
Only the original pointwise determinant hypothesis 1 is required.
\end{theorem}
\begin{proof}
Again suppose all lengths finite. Choose lifts in $R$ of a basis
of the joint residual image $S$. Parameterize each entry of
$C\in M_{b\times a}(R)$ by $O$-linear combinations of these lifts.
Fix $O$-bases of $W_j$ and put
\[
                 \Phi(C)=\prod_j\det(\rho_j(C)|_{W_j}).
\]
The reduction of this polynomial is the nonzero product in
Proposition~\ref{prop:residual-block-large-field}. Thus $\Phi$ has
a unit coefficient and total degree $D_0$.
Lemma~\ref{lem:residual-block-grid} gives $C$ over $R$ with
$v(\Phi(C))\le v(D_0!)$. Each restricted matrix is invertible over
$K$. Relative to a basis of $W_j$, elementary divisors give
\[
 v(\det\rho_j(CP))
 =L_j+v(\det(\rho_j(C)|_{W_j})).
\]
The determinant function of $CP$ still belongs to $T$ and
vanishes away from $x_0$. Summing this identity proves
\eqref{eq:residual-block-factorial}. The normalized formulas follow
from $v(n)=e_Ov_\ell(n)$ for nonzero integers and the factorization
of a factorial into prime powers.
\end{proof}

\begin{corollary}[The total factorial correction]
\label{cor:residual-block-factorial-sum}
Let $0\le D_\ell\le D$ be integers and $0\le w_\ell\le W$.
Then
\[
 \sum_\ell w_\ell v_\ell(D_\ell!)\log\ell
       \le W\log(D!)
       \le WD\log(\max(2,D)).
\]
For several applications indexed by $a$, the same bound holds if
$D_{a,\ell}\le D$ and $\sum_a w_{a,\ell}\le W$ for every $\ell$.
If there are at most $m$ applications, each of weight at most $W$,
the bound is $mW\log(D!)$.
\end{corollary}
\begin{proof}
Use $v_\ell(D_\ell!)\le v_\ell(D!)$ and
$\sum_\ell v_\ell(D!)\log\ell=\log(D!)$. The same calculation
uses the total weight at each prime in the multiple-application case.
\end{proof}

\begin{proposition}[Normalized coefficient extension]
\label{prop:residual-block-base-change}
Let $O'/O$ be a finite extension of DVRs with relative ramification
index $h$. For a finite-length $O$-module $C$,
\[
               \operatorname{length}_{O'}(C\otimes_OO')
                   =h\operatorname{length}_O C.
\]
Supported congruence ideals commute with this extension. For a
lattice $M$ in a decomposed $K$-vector space and projection $\pi_Y$,
the quotient $\pi_Y(M)/(M\cap V_Y)$ also commutes with it.
In mixed characteristic with residue characteristic $\ell$, the
lengths of these quotients divided by $v_O(\ell)$ are unchanged.
\end{proposition}
\begin{proof}
The ring $O'$ is flat over $O$. Elementary divisors reduce the
length assertion to $C=O/(\varpi^a)$, whose scalar extension has
length $ha$. The submodule of $T$ supported at $x_0$ is the kernel
of $T\to O^{X\setminus\{x_0\}}$; flatness preserves this kernel
and its image under evaluation at $x_0$. It likewise preserves the
kernels and images defining intersections and projections of the
lattice, and hence their quotient. Finally
$v_{O'}(\ell)=h v_O(\ell)$, so the two factors cancel.
\end{proof}

\begin{proposition}[Exact bound after enlarging the residue field]
\label{prop:residual-block-stable-law}
Suppose $O$ is complete, each $\bar\rho_j$ is onto, and assumption 1
remains valid after every finite unramified extension: for each such
$O'/O$ and every $Q'\in M_b(R\otimes_OO')$, its determinant tuple
belongs to $T\otimes_OO'$. Then
\eqref{eq:residual-block-exact} holds at every residue characteristic.
\end{proposition}
\begin{proof}
If the residue field is infinite or already has size greater than
$D_0$, apply Proposition~\ref{prop:residual-block-large-field}
directly. Otherwise it is finite, and completeness supplies a
finite unramified extension with residue-field size greater than
$D_0$. Individual residual surjectivity persists
after base change. Apply
Proposition~\ref{prop:residual-block-large-field} and descend the
inequality by Proposition~\ref{prop:residual-block-base-change},
with relative ramification index one.
\end{proof}

The stable determinant hypothesis is an additional requirement.
Pointwise membership over $O$ alone need not imply membership
after extension: a polynomial such as $Z^q-Z$ vanishes on every
point of the residue field of order $q$ while being a nonzero
polynomial. One must verify a determinant polynomial law or its
concrete arithmetic construction before using
Proposition~\ref{prop:residual-block-stable-law}.

Nor can an unramified extension always split the characteristic-zero
eigensystems: $O[T]/(T^2-\varpi)$ needs a ramified extension to
contain $\sqrt\varpi$. Proposition~\ref{prop:residual-block-base-change}
permits that extension without changing normalized defects.

\paragraph{Applicability to cyclotomic blocks.}
The residual characters $\bar\chi^a$ and $\bar\chi^{a+\ell-1}$
coincide. Their joint image is diagonal, so the source's stronger
joint-surjectivity hypothesis cannot be checked for a long
consecutive family literally as written. Each scalar block is
individually onto, because $R$ is an $O$-algebra. Thus the preceding
individual-image variants give a corrected algebraic interface.
The off-target kernels, full residual image of any nonscalar block,
determinant law when needed, local Selmer estimates, and the actual
number and weights of applications remain separate arithmetic
hypotheses. The factorial error is polynomial up to a logarithm
when the matrix dimensions, weights and application count are
polynomial in the scale; its absorption into the required global
budget still has to be checked there.
